\documentclass[11pt,a4paper]{article}
\usepackage[margin=1.1in]{geometry}
\usepackage{amsmath,amssymb,amsthm,mathtools}
\usepackage[T1]{fontenc}
\usepackage{microtype}
\usepackage{enumitem}
\usepackage{booktabs}
\usepackage[hidelinks]{hyperref}

\newtheorem{theorem}{Theorem}[section]
\newtheorem{proposition}[theorem]{Proposition}
\newtheorem{lemma}[theorem]{Lemma}
\newtheorem{corollary}[theorem]{Corollary}
\theoremstyle{definition}
\newtheorem{definition}[theorem]{Definition}
\newtheorem{remark}[theorem]{Remark}
\theoremstyle{remark}
\newtheorem{scope}[theorem]{Scope}
\newtheorem{escape}{Escape route}

\newcommand{\Z}{\mathbb{Z}}
\newcommand{\Q}{\mathbb{Q}}
\newcommand{\R}{\mathbb{R}}
\newcommand{\HH}{\mathbb{H}}
\newcommand{\vol}{\operatorname{vol}}
\newcommand{\Hom}{\operatorname{Hom}}

\title{What a class invariant cannot supply:\\
scale, orbit position, and the difference between\\
a manifold and its family}
\author{}
\date{26 August 2026}

\begin{document}
\maketitle

\begin{abstract}
Programmes that extract structure from a fixed hyperbolic $3$-manifold routinely reach a point
where a desired quantity fails to appear, and it is rarely clear whether the failure is a gap in
the argument or a theorem. We isolate three such failures for the figure-eight knot complement
$m004$ and show that each is a theorem, with an exactly stated escape route.

\emph{Scale.} Mostow rigidity fixes the shape of a finite-volume hyperbolic $3$-manifold but not
its size, and the covering ladder satisfies $\vol(\widetilde X)=d\cdot\vol(X)$ with no
distinguished degree. Consequently no dimensionful quantity is determined. The escape is exact
and is not a loophole: dimensionless quantities --- ratios, phases, mixing angles, counts --- are
untouched by the argument, so the no-go constrains what kind of quantity may be sought, not
whether any may be.

\emph{Orbit position.} If a quantity is an invariant of an orbit of a group acting on a space of
candidates, it is constant on that orbit and therefore cannot select a point of it. This is
structural, not evidential: no refinement of the same invariant will ever select a point, so
accumulating more of the same evidence is provably futile. There are exactly two escapes ---
shrink the group, or add structure that is not invariant.

\emph{Family versus object.} Of the fourteen orientable cusped census manifolds whose tetrahedron
shape field is $\Q(\sqrt{-3})$, we tabulate seven elementary invariants and find that
\emph{exactly one} separates $m004$ from the rest: $H_1\cong\Z$. Volume, tetrahedron count, cusp
count, amphichirality and vanishing Chern--Simons invariant are all shared. In particular the
identity $\vol(m004)=\tfrac{3\sqrt3}{2}L(\chi_{-3},2)$ is a property of the family, being shared
with $m003$. It follows that any construction taking the invariant trace field as its input has a
\emph{family-level} input, and cannot by itself distinguish $m004$ from thirteen other manifolds.

We report two errors of our own that this census caught, because both are instructive: comparing
Chern--Simons invariants by floating-point equality, which produced a spurious separating
property; and testing torsion-freeness in place of the knot-complement condition, which is a
different property that two other family members also satisfy.
\end{abstract}

\section{Introduction}

\subsection{Negative results as instruments}

A programme that derives structure from a fixed mathematical object accumulates, alongside its
positive results, a list of things it has failed to obtain. That list is usually treated as a
record of unfinished business. Sometimes it is instead a record of theorems, and the difference
matters: an unfinished computation rewards more effort, while a theorem of impossibility tells
one that a particular kind of effort is wasted and which other kind is not.

This paper takes three items from such a list and determines which they are. All three turn out
to be theorems. In each case we state the theorem, then state the escape route precisely, because
a no-go result without its escape route is misleading --- it invites the reader to conclude that
more is closed than actually is.

\subsection{The object}

Throughout, $X$ denotes a finite-volume orientable hyperbolic $3$-manifold, and $m004$ the
figure-eight knot complement, with $\vol(m004)=2.029883212819307\ldots$ and invariant trace field
$\Q(\sqrt{-3})$.

\section{Scale is not determined}\label{sec:scale}

\begin{theorem}[no dimensionful quantity]\label{thm:scale}
Let $F$ be a function assigning a real number to a finite-volume hyperbolic $3$-manifold, defined
purely in terms of its isometry type. Then $F$ cannot determine a quantity carrying a dimension of
length, area or volume in any way that distinguishes a manifold from its finite covers, because
for every degree-$d$ cover $\widetilde X\to X$ one has $\vol(\widetilde X)=d\cdot\vol(X)$, and the
covering ladder above $X$ contains covers of unboundedly many degrees with no distinguished one.
\end{theorem}

\begin{proof}
Mostow rigidity makes the hyperbolic structure on $X$ unique, so isometry-type invariants are
topological invariants; they do not depend on a choice of unit. Volume multiplies exactly by the
covering degree, and a finite-volume hyperbolic $3$-manifold has finite covers of infinitely many
degrees. Any rule selecting one member of the ladder is additional data, not a consequence of the
isometry type.
\end{proof}

\begin{escape}[dimensionless quantities are untouched]\label{esc:ratios}
Theorem~\ref{thm:scale} constrains dimensionful quantities only. A ratio of two quantities that
scale the same way, a phase, a mixing angle, or an integer count is invariant under the ladder and
is not excluded by any part of the argument. We emphasise this because the theorem is easy to
overstate into ``nothing numerical can be obtained'', which is false and would close off the one
lane the argument leaves open.
\end{escape}

\section{An orbit invariant cannot select a point of the orbit}\label{sec:orbit}

\begin{theorem}[the orbit-to-point gap is structural]\label{thm:orbit}
Let a group $G$ act on a set $\Omega$ of candidate objects, and let $I:\Omega\to V$ be
$G$-invariant. Then $I$ is constant on each $G$-orbit, and therefore no function of $I$ alone
distinguishes two points of the same orbit. In particular, if a desired conclusion amounts to
selecting a point of a $G$-orbit, no refinement of $I$ --- no sharper computation of the same
invariant, no additional decimal places, no further consistency check --- can ever supply it.
\end{theorem}

The proof is immediate. We state it as a theorem because its content is not the mathematics but
the classification it forces: a failure to select a point of an orbit is not weak evidence, it is
a proof that the evidence is of the wrong type. Effort spent sharpening $I$ is provably wasted.

\begin{escape}[there are exactly two]\label{esc:orbit}
To select a point of a $G$-orbit one must either
\begin{enumerate}[label=\textup{(\roman*)},leftmargin=2.4em]
\item \textbf{shrink the group}: replace $G$ by a subgroup $H$ for which the orbit breaks up, so
      that the $H$-orbits are smaller and $I$ becomes a finer invariant --- for instance replacing
      a $\Q$-linear symmetry by a $\Z$-linear one, where the failure of homogeneity is exactly what
      creates new invariants; or
\item \textbf{add non-invariant structure}: introduce data not preserved by $G$, which is to say,
      change the problem to one where the point is not merely a point of an orbit.
\end{enumerate}
These are exhaustive, since any successful selector either fails to be $G$-invariant (case (ii))
or is invariant for a smaller group (case (i)).
\end{escape}

\begin{remark}
Escape~\ref{esc:orbit} is useful precisely because it is short. It converts an open-ended search
into a choice between two named moves, each of which can be attempted or ruled out on its own
terms.
\end{remark}

\section{The family, and what separates $m004$ from it}\label{sec:family}

The invariant trace field is a strong invariant, but it is an invariant of a \emph{family}. We
make this quantitative.

\begin{definition}
Let $\mathcal{F}$ be the set of orientable cusped census manifolds whose tetrahedron shape field is
$\Q(\sqrt{-3})$. Then $|\mathcal{F}|=14$:
\[
\begin{aligned}
 \mathcal{F}=\{\,&m003,\ m004,\ m202,\ m203,\ m206,\ m207,\ m208,\\
                 &m410,\ m412,\ s118,\ s119,\ s594,\ s595,\ s596\,\}.
\end{aligned}
\]
\end{definition}

\begin{theorem}[exactly one elementary invariant separates $m004$]\label{thm:family}
Among the seven invariants --- volume, tetrahedron count, cusp count, torsion-freeness of $H_1$,
the condition $H_1\cong\Z$, amphichirality, and vanishing of the Chern--Simons invariant --- exactly
one distinguishes $m004$ within $\mathcal{F}$, namely $H_1\cong\Z$. The other six are shared:
\begin{center}
\begin{tabular}{@{}lll@{}}
\toprule
invariant & value at $m004$ & also held by \\
\midrule
volume                   & $2.029883212819307$ & $m003$ \\
tetrahedron count        & $2$                 & $m003$ \\
cusp count               & $1$                 & nine others \\
torsion-free $H_1$       & true                & $m202$, $m203$ \\
amphichiral              & true                & \textbf{all thirteen others} \\
Chern--Simons $=0$       & true                & $m203$, $m206$, $m208$, $s595$, $s596$ \\
\midrule
$H_1\cong\Z$             & true                & \textbf{none} \\
\bottomrule
\end{tabular}
\end{center}
\end{theorem}

\begin{corollary}[the volume identity is a family property]\label{cor:volume}
The identity $\vol(m004)=\tfrac{3\sqrt3}{2}L(\chi_{-3},2)$ is shared with $m003$ and is therefore
a property of the family, not of $m004$. It cannot participate in any argument that needs to
distinguish $m004$ from other members of $\mathcal{F}$.
\end{corollary}

\begin{corollary}[trace-field inputs are family inputs]\label{cor:familyinput}
Any construction whose input is the invariant trace field takes a family-level input: all fourteen
members of $\mathcal{F}$ supply the same field and would proceed identically. Such a construction
cannot, by itself, be said to be about $m004$ rather than about $\mathcal{F}$.
\end{corollary}

\begin{scope}[what Corollary~\ref{cor:familyinput} does and does not damage]\label{sc:notdamaging}
This is not a contradiction of arguments that first select $m004$ and then use the trace field: if
the selection step uses $H_1\cong\Z$, it is genuinely object-level by
Theorem~\ref{thm:family}, and a later trace-field step inherits that selection. What
Corollary~\ref{cor:familyinput} forbids is the reverse order, and it sharpens the honest
description of any such chain: the field-based step is a family step, and should be described as
one. A reader who runs this census will notice, and it is better to state it first.
\end{scope}

\section{Two errors this census caught}\label{sec:errors}

Both errors were ours, both produced confident wrong answers, and both were caught by the census
rather than by inspection. We report them because each is a general failure mode.

\begin{enumerate}[leftmargin=1.8em]
\item \textbf{Floating-point equality on an invariant that is only numerically zero.} Our first
      run compared Chern--Simons invariants with exact equality. The computed value at $m004$ was
      $9\times10^{-17}$ rather than $0$, so $m004$ appeared to be the unique member with vanishing
      Chern--Simons invariant, and ``$CS=0$'' was reported as a second separating property. It is
      not: five other members have $CS=0$ as well. The error inverted the paper's conclusion in
      the direction that flattered it.
\item \textbf{Testing a neighbouring property.} Our first run tested torsion-freeness of $H_1$,
      which is \emph{not} the knot-complement condition. The manifolds $m202$ and $m203$ have
      $H_1\cong\Z\oplus\Z$, which is torsion-free but not $\Z$. Substituting a property that is
      easier to compute for the one actually at issue is a failure mode that no amount of care
      \emph{within} the computation detects, because the computation is correct --- it is
      answering a different question.
\end{enumerate}

The first was caught by re-reading the run's own printed table rather than its summary line; the
second by writing down the intended property and the tested property side by side.

\section{Summary}

\begin{center}
\begin{tabular}{@{}p{4.5cm}lp{5.4cm}@{}}
\toprule
failure & status & escape \\
\midrule
no dimensionful quantity & Thm~\ref{thm:scale} & dimensionless quantities untouched (\ref{esc:ratios}) \\
cannot select within an orbit & Thm~\ref{thm:orbit} & shrink the group, or add non-invariant structure (\ref{esc:orbit}) \\
trace field does not isolate $m004$ & Thm~\ref{thm:family} & select first by $H_1\cong\Z$ (\ref{sc:notdamaging}) \\
\bottomrule
\end{tabular}
\end{center}

Each is closed as a matter of evidence and open as a matter of route. That combination is the
useful output: it says which further computation is futile, and names the alternatives that are
not.

\appendix

\section{Verification}\label{app:verify}

The census of \S\ref{sec:family} is reproduced by \texttt{verify/}, which rebuilds $\mathcal{F}$
from the census by shape field rather than accepting a stored list, and tabulates all seven
invariants. Controls: Chern--Simons is compared with an explicit tolerance and never by equality,
with the tolerance stated; the intended property ($H_1\cong\Z$) and the neighbouring property
(torsion-freeness) are both computed and reported separately, so that substituting one for the
other is visible rather than silent; and the family is regenerated, so a change in the census
would change the table rather than pass unnoticed.

\section{Status in the programme register}\label{app:register}

\begin{center}
\begin{tabular}{@{}llp{6.9cm}@{}}
\toprule
entry & status & relation to this paper \\
\midrule
\texttt{OA-C0016} & \texttt{EXTERNAL\_BLOCKER} & whether every dimensionless parameter is emitted after thresholds and running. Untouched here; \S\ref{sec:scale} constrains only dimensionful quantities. \\
\texttt{OA-C0018} & \texttt{EMPIRICAL} & whether the completed theory makes a successful unused quantitative prediction. Not addressed. \\
\texttt{OA-C1063} & \texttt{REFUTED} & whether a Beilinson regulator is defined for the relevant classes. Consistent with the orbit argument of \S\ref{sec:orbit}. \\
\texttt{OA-C1064} & \texttt{EXTERNAL\_BLOCKER} & whether an arithmetic motive is uniquely selected. \S\ref{sec:family} explains one reason selection is hard: the natural inputs are family-level. \\
\bottomrule
\end{tabular}
\end{center}

\noindent
None of these is \texttt{PROVED}. This paper does not change that, and does not try to: its
contribution is to convert three of the surrounding failures from open questions into theorems
with named escape routes, which is a different and smaller thing than answering them.

\begin{thebibliography}{9}
\bibitem{Mostow} G.~D.~Mostow, \emph{Strong Rigidity of Locally Symmetric Spaces}, Princeton, 1973.
\bibitem{Prasad} G.~Prasad, Strong rigidity of $\Q$-rank $1$ lattices, \emph{Invent.\ Math.}
  \textbf{21} (1973), 255--286.
\bibitem{MaclachlanReid} C.~Maclachlan and A.~W.~Reid, \emph{The Arithmetic of Hyperbolic
  $3$-Manifolds}, Springer, 2003.
\bibitem{Thurston} W.~P.~Thurston, \emph{The Geometry and Topology of Three-Manifolds},
  Princeton lecture notes, 1980.
\bibitem{CullerDunfield} M.~Culler, N.~M.~Dunfield, M.~Goerner and J.~R.~Weeks, SnapPy, a computer
  program for studying the geometry and topology of $3$-manifolds,
  \texttt{http://snappy.computop.org}.
\bibitem{NeumannZagier} W.~D.~Neumann and D.~Zagier, Volumes of hyperbolic three-manifolds,
  \emph{Topology} \textbf{24} (1985), 307--332.
\end{thebibliography}

\end{document}
